\chapter*{پیشگفتار }

"نظریهٔ انواع" یکی از مهم‌ترین شاخه‌های علوم کامپیوتر نظری و منطق ریاضی در قرن بیستم به شمار می‌رود. 
این نظریه نه تنها پایه و اساس زبان‌های برنامه‌نویسی تابعی پیشرفته مانند
 \lr{Haskell}، 
 \lr{OCaml} 
و
 \lr{Agda} 
 را تشکیل می‌دهد، بلکه هسته اصلی اثبات‌یارهای مطرحی نظیر
 \lr{Coq} 
و
 \lr{Lean}
 نیز بر مبنای آن طراحی شده است.

یکی از مسائل کلاسیک و همچنان باز در این حوزه، رابطهٔ میان نرمال‌سازی ضعیف و نرمال‌سازی قوی در سیستم‌های نوع وابسته است.
 حدس معروف
 \lr{BGK} 
که بیش از سه دهه پیش مطرح شد، بیان می‌کند که 
در هر سیستم نوع محض، اگر نرمال‌سازی ضعیف برقرار باشد، آنگاه نرمال‌سازی قوی نیز برقرار خواهد بود.

این حدس برای سیستم‌های غیروابسته از دهه ۱۹۹۰ اثبات شده است
 \cite{barthe2001}
، اما در سیستم‌های وابسته مانند حساب ساختمان‌ها
 \footnote{
    هسته اثبات‌یار
     \lr{Coq}
 }
 همچنان یک مسئله باز باقی مانده است.

\lr{Nathan Mull} 
در سال ۲۰۲۳، در رساله دکتری خود در دانشگاه شیکاگو
 \cite{mull2023}
، خانواده جدیدی از سیستم‌های نوع به نام «سیستم‌های نوع محض لایه‌ای
\footnote{\lr{Tiered Pure Type Systems}}» 
را معرفی کرد و با ارائهٔ ترجمه‌ای هوشمندانه که وابستگی‌های غیرضروری را حذف می‌کند، 
توانست نشان دهد که در این سیستم‌ها، حدس
 \lr{BGK} 
برقرار است. 
این نتیجه یکی از قوی‌ترین پیشرفت‌های اخیر در مسیر اثبات این حدس به‌ شمار می‌رود.

این پایان‌نامه به صوری‌سازی کامل اثبات
 \lr{Mull} 
در اثبات‌یار
 \lr{Coq} 
اختصاص دارد. 
هدف، تبدیل این اثبات نظری قابل توجه، به یک اثبات ماشینی قابل اتکا و استفاده مجدد است که 
می‌توان از آن برای تعمیم به سیستم‌های پیچیده‌تر نیز استفاده کرد.

این پایان‌نامه در پنج فصل تنظیم شده است. 
فصل اول مفاهیم و تعاریف اولیه را بیان می‌کند. 
فصل دوم سیستم‌های نوع محض لایه‌ای و ترجمه حذف وابستگی در آنان را معرفی می‌کند. 
فصل سوم به طراحی و پیاده‌سازی سیستم‌ها و توابع ترجمه در
 \lr{Coq} 
می‌پردازد. 
فصل چهارم اثبات‌های صوری لم‌ها و قضیه اصلی را در
 \lr{Coq} 
ارائه می‌دهد 
و فصل پنجم به جمع‌بندی و پیشنهادهایی برای پژوهش‌های آینده اختصاص دارد.

امید است این کار بتواند گامی کوچک در جهت پیشرفت یکی از مسائل کلاسیک نظریه انواع بردارد.